% !TEX encoding = UTF-8 Unicode

% Beispiel für ein LaTeX-Dokument im Format "seminarvorlage"
\documentclass[ngerman]{article}

\usepackage[utf8]{inputenc}
\usepackage{csquotes}
\usepackage{float}
\usepackage{graphicx}
\usepackage{hyperref}
\usepackage[acronym,nonumberlist]{glossaries-extra}
\usepackage{cleveref}
\usepackage{nameref}
\usepackage{array}
\usepackage{multirow}
\usepackage{listings}
\usepackage{parskip}
\usepackage{attachfile2}
\usepackage{fancyhdr}
\usepackage{ragged2e}
\usepackage{multicol}
\usepackage{geometry}
\usepackage{changepage}
\usepackage{enumitem}
\usepackage{tablefootnote}
\usepackage{threeparttable}
% ngerman = Deutsch in neuer Rechtschreibung, alternativ english
\usepackage{babel} % automatische Sprachunterstützung

\usepackage[utf8]{inputenc} % Kodierung der Non-ASCII-Zeichen
\usepackage[T1]{fontenc} % Moderne Fonts, Trennung von Wörtern mit Umlauten
\usepackage{cleveref} % für bequeme Referenzen, siehe \cref unten

\usepackage{embrac}% upright brackets in emphasised text, () and [], empfohlen

%\usepackage{biblatex}% besser als bibtex, aber dann biber statt bibtex benutzen

\begin{document}

\title{Evaluation von KI-Funktionen in modernen Entwicklungsumgebungen}
\author{Konstantin Fastovski, Silas Rißler}
\begin{titlepage}
    \begin{center}
        \vspace{1cm}

        \Huge
        Evaluation von KI-Funktionen in modernen Entwicklungsumgebungen

        \LARGE
        \vspace{1.5cm}
        \textbf{Silas Rißler}

        \textbf{Konstantin Fastovski}

        \vfill

        \Large
        Angewandte Informatik \\
        Duale Hochschule Baden-Württemberg Stuttgart Campus Horb \\
        \vspace{1.5cm}

        Oktober 2026

        \vspace{1cm}
    \end{center}
\end{titlepage}

\newpage

\section{Exposé}

Dieser Artikel dient der Gegenüberstellung mehrerer beliebter Entwicklungsumgebungen bzw. Text-Editoren im Bezug auf die Nutzerfreundlichkeit, Transparenz und Qualität der von ihnen angebotenen Werkzeuge zur Einbindung von künstlicher Intelligenz (KI) in den Workflow von Software-Entwicklern. Es wird identifiziert, welche Design-Prinzipien sich bei mehreren Produkten durchgesetzt haben und weshalb. Weiterhin wird evaluiert, welche Fehler unbeliebte Umsetzungen von KI-Einbindungen machen.

Hierzu werden einige Code-Editoren, Integrated Development Environments (IDEs) und andere Programme getestet und verglichen, die KI in den Prozess der Software-Entwicklung einbeziehen. Es wird analysiert, welche Entscheidungen in der Gestaltung der Features Nutzer*innen das Vertrauen in die Ausgabe nehmen oder intransparent mit dem Eingabekontext umgehen, was zwei der Hauptgründe von Entwickler*innen für das Nicht-Benutzen von KI bei ihrer Arbeit sind. \cite{Sergeyuk_2025}

Es wird eine Datenlage geschaffen, durch die die wesentlichen Eigenschaften der Implementierungen von beispielsweise KI-Chatfenstern deutlich werden, die für Nutzer*innen von Bedeutung sind. Weiterhin werden Features nach bestimmten Anwendungsfällen --- wie zum Beispiel das Erkunden von Lösungen gegenüber konkreter Implementierung --- gruppiert und deren Qualität in bestimmten Einsatzgebieten untersucht, um mögliche Parallelen zu der statistischen Verteilung von Aufgabentypen zu erkennen, für die Entwickler*innen häufiger bzw. seltener zu KI greifen. \cite{barke2022groundedcopilotprogrammersinteract}

\bibliographystyle{abbrv}
\bibliography{literatur} % Daten aus der Datei literatur.bib verwenden.

\end{document}
